% Abstract (ingles).
\begin{otherlanguage}{english}
\chapter*{Abstract}
\addcontentsline{toc}{chapter}{Abstract}
\markboth{Abstract}{Abstract}

This thesis presents a reproducible urban growth model for the Querétaro
Metropolitan Area, built from a 37-year annual image series (1984--2020) and
the coupling of Weights of Evidence (WoE) with cellular automata (CA). RGB
captures obtained through Google Earth's historical imagery tool were turned
into binary urban/non-urban maps via a classification pipeline based on local
binary patterns, K-Means clustering and a linear SVM classifier. On these maps,
the Weights of Evidence quantify the spatial influence of a small set of
auditable variables and feed, together with the Moore neighborhood, a
probabilistic transition rule that the cellular automaton applies cell by cell.

The model was calibrated through a systematic search over the transition
threshold, the neighborhood weight and the \textit{Information Value} weighting,
and was assessed with a validation protocol across five independent five-year
windows (2011--2016 to 2015--2020). Average performance reaches a Figure of
Merit of $0.317$, a Kappa coefficient of $0.445$, an accuracy of $0.722$ and an
IoU of $0.633$, using the Pontius error decomposition as the interpretation
framework. The central contribution is not the WoE-CA coupling itself, but a
reproducible multi-temporal validation protocol, applied to a scarcely
documented Mexican context, with a white-box implementation that runs entirely
on CPU, without GPU or proprietary software, from openly accessible data.

\vspace{0.6cm}
\noindent\textbf{Keywords:} urban growth; cellular automata; Weights of
Evidence; multi-temporal validation; Figure of Merit; Querétaro Metropolitan
Area.
\end{otherlanguage}
\cleardoublepage
